\section{输运结构与局部生成映射}
\label{sec:theory}

力学状态采用归一化向量 $z=(Q,P)\in U\subset\mathbb R^{2n}$，时间为 $\tau$，Hamilton 函数为 $h$，步长为 $\eta=\Delta\tau$。杆降阶模型中 $n=12$，但以下构造与这一维数无关。这里 $M$ 表示归一化质量矩阵。工作域取为质量矩阵正定且所述导数存在的开集。所有局部流、辅助流及隐式分支都必须留在指定域内。约定
\begin{equation}
 J=\begin{pmatrix}0&I\\-I&0\end{pmatrix},\qquad K_0=-J,
 \qquad (Da)_{ij}=\partial_{z_j}a_i .
 \label{eq:th-conventions}
\end{equation}
这一约定统一了变分方程和端点构造的符号。几何积分的基本框架见 \citet{FengQin2010,Hairer2006}；这里明确非可分力学模型所需的相容性与可解性条件。

\subsection{具有可测残差的变分表示}
% Source: stage5/proofs/corrected_mathematics.md, section 3, R7--R10; stage4/audit/birkhoff_full_audit.md, section 2; original1025 p1113--p1125, p1180.
令 $\alpha_\tau=a(z,\tau)^T dz$，$\Omega_\tau=d\alpha_\tau$，则
\begin{equation}
 \Omega_\tau=\frac12\sum_{i,j}K_{ij}\,dz_i\wedge dz_j,
 \qquad K=(Da)^T-Da .
 \label{eq:th-form}
\end{equation}
对 $\int(a^T\dot z-B)\,d\tau$ 作固定端点变分并分部积分，得到
\begin{equation}
 \begin{aligned}
 \delta\mathcal A&=\int\delta z^T(K\dot z-\partial_\tau a-\nabla B)\,d\tau,\\
 Kf&=\nabla B+\partial_\tau a .
\end{aligned}
 \label{eq:th-variation}
\end{equation}
因此若引入 $D=-\nabla B-\partial_\tau a$，向量场应为 $-K^{-1}D$。闭性由 $d^2=0$ 得到；$K$ 的可逆性是独立假设。Cartan 恒等式和 $\iota_f\Omega=-(Kf)^Tdz$ 给出
\begin{equation}
 \begin{aligned}
 \partial_\tau\Omega+\mathcal L_f\Omega&=0,\\
 K_\tau+DK[f]+Df^TK+KDf&=0.
\end{aligned}
 \label{eq:th-compatibility}
\end{equation}
于是流满足 $\Phi_{\tau,s}^*\Omega_\tau=\Omega_s$。所保持的是两个时间截面之间的体积 $\Omega_\tau^n/n!$。它通常不同于欧氏相体积：统一动量阻尼使后者收缩，同时保持一个指数加权二形式。

% Source: stage4/audit/birkhoff_full_audit.md, section 3, recurrence (T); stage5/proofs/birkhoff_transport.md, sections 1--2 and 6; stage5/program/transport_native.py, lie_oneform and source_coefficient_function.
对于光滑且局部可逆的流，$\alpha_\tau=(\Phi_{\tau,s}^{-1})^*\alpha_s$ 给出局部存在性构造。相应微分方程与标量为
\begin{equation}
 \begin{aligned}
 \partial_\tau a+\mathcal L_fa&=0,\qquad B=a^Tf,\\
 \mathcal L_fa&=Da\,f+Df^Ta.
\end{aligned}
 \label{eq:th-homogeneous}
\end{equation}
计算有限时间展开不需要预先知道物理精确流。令 $\theta=\tau-s$，写成 $f=\sum_{j\ge0}\theta^j f_j$ 和 $a=\sum_{j\ge0}\theta^j a_j$，其中系数已包含阶乘分母。逐幂比较得到
\begin{equation}
 a_{m+1}=-\frac1{m+1}\sum_{j+k=m}
       \bigl(Da_j f_k+Df_k^Ta_j\bigr).
 \label{eq:th-time-recursion}
\end{equation}
显式时间依赖必须进入 $f_k$。有限个系数不能证明无限级数收敛。

对实际实现的自治力学分支，写成
$h=T+V$，$T=P^TM(Q)^{-1}P/2$，以及 $f=J\nabla h+(0,r)$。取 $a_0=(P,0)$、$\ell=P^Th_P-h=T-V$，直接求导得 $\mathcal L_fa_0=d\ell+(r,0)$。对一个等价的有源方程作截断更为合适：
\begin{equation}
 \begin{aligned}
 a'&=a_0+b,\qquad B'=h+b^Tf,\\
 \partial_\tau a'+\mathcal L_fa'&=d\ell,\\
 b_1&=-(r,0),\\
 b_{j+1}&=-\frac{\mathcal L_fb_j}{j+1}.
\end{aligned}
 \label{eq:th-source}
\end{equation}
这里 $b^{[r_o]}=\sum_{j=1}^{r_o}\theta^jb_j$ 的截断阶数为 $r_o$，与载荷 $r$ 区分。其表示残差严格为
\begin{equation}
 \begin{aligned}
 e_B&=K^{[r_o]}f-\partial_\tau a'^{[r_o]}-\nabla B'^{[r_o]}\\
    &=-\theta^{r_o}\mathcal L_fb_{r_o},\\
 f_a&=f-(K^{[r_o]})^{-1}e_B .
\end{aligned}
 \label{eq:th-source-residual}
\end{equation}
当载荷为零时，所有 $b_j$ 都为零，每个截断阶数均恢复原 Hamilton 系统。标量方程 $\sigma_\tau+f^T\nabla\sigma=\ell$ 通过 $a=a'-d\sigma$、$B=B'+\sigma_\tau$ 将此构造与齐次输运联系起来。规范变换保持表示残差；截断与规范变换一般不可交换。因此这里通过直接截断有源方程本身来改善有限阶近似。

% Source: stage5/proofs/birkhoff_transport.md, sections 3 and 7; stage5/program/transport_cache.py, OneFormCache and CachedRepresentedField; stage5/proofs/high_derivative_ad.md, replacement-path discussion.
数值坐标图在固定中心使用这些一形式系数的空间 Taylor 多项式 $\widetilde b_j$，保留 $h$、$f$、$Df$ 的完整求值器。然而，被表示的场仍发生变化：取 $\widetilde B=h+\widetilde b^Tf$ 和 $\widetilde K=K_0+D\widetilde b^T-D\widetilde b$，则
\begin{equation}
 \begin{aligned}
 \widetilde e_B&=-\bigl(\widetilde b_\tau+D\widetilde b f
                         +Df^T\widetilde b+(r,0)\bigr),\\
 \widetilde f_a&=f-\widetilde K^{-1}\widetilde e_B .
\end{aligned}
 \label{eq:th-cache-residual}
\end{equation}
因此时间截断、系数的空间近似与数值积分各有独立误差。高阶导数使用代数等价表达式 $P^TM(Q)^{-1}P/2+h(Q,0)$ 求值，质量矩阵在当前位形重新计算。独立的反向模式方向导数和有限差分检验所选求导路径；仅有代数等价或张量对称性，不能验证自动微分实现正确。

独立导数检验的必要性可由标量函数 $h(Q,P)=P^2/[2(1+Q^2)]$ 看出。其展开给出 $h_{QQPP}(0)=h_{PPQQ}(0)=-2$。在所检验的软件版本中，一条经过线性求解的混合求导路径将一个四阶分量计算为零。张量指标置换平均一般不能纠正这种错误，因为与同一方向的重复副本收缩时，平均不改变结果。因此，对替代表达式还须数值交叉检查输运系数实际使用的导数量。这些检查提供的是针对所选路径、点和方向的证据。

这一区分直接影响计算解释：即使被表示的向量场与目标场存在明显差异，坐标图积分器仍可准确保持缓存的闭二形式。在缓存中心，精确的系数值和一阶导数可重现原生时间截断残差。离开中心后，须独立计算原生系数才能评估空间近似。小的拉回缺陷或隐式残差均不能界定这一差异。当前构造使用一个固定时间原点和空间域；若通过重叠局部坐标图延拓，还须保证过渡映射和二形式一致。

\subsection{可计算的正则坐标图及其时间导数}
% Source: stage5/proofs/birkhoff_transport.md, section 4; stage5/program/moser_chart.py, auxiliary, map, transformed_field, transformed_jacobian.
在每个固定物理时间，令
\begin{equation}
 \begin{aligned}
 K_{\lambda,\tau}&=K_0+\lambda(\widetilde K_\tau-K_0),\\
 v_{\lambda,\tau}&=K_{\lambda,\tau}^{-1}\widetilde b(\cdot,\tau),\\
 \partial_\lambda\rho_{\lambda,\tau}
     &=v_{\lambda,\tau}\circ\rho_{\lambda,\tau}.
\end{aligned}
 \label{eq:th-moser}
\end{equation}
其中 $\rho_{0,\tau}=\mathrm{id}$。$\widetilde K_\tau$ 的下标 $\tau$ 表示时间截面。此同伦要求每个 $\lambda\in[0,1]$ 都非退化，仅有端点非退化还不够。因为 $\iota_v\Omega_{\lambda,\tau}=-\widetilde b^Tdz$，对拉回求导得 $\rho_{1,\tau}^*\widetilde\Omega_\tau=\Omega_0$。因此 $\chi_\tau=\rho_{1,\tau}^{-1}$ 是正则坐标图。

记 $\rho=\rho_{1,\tau}$、$\mathsf R=D_y\rho$、$\upsilon=\partial_\tau\rho$。正则坐标中的向量场为
\begin{equation}
 u(y,\tau)=\mathsf R^{-1}\{\widetilde f_a(\rho(y,\tau),\tau)-\upsilon\}.
 \label{eq:th-chart-field}
\end{equation}
其中包含坐标图的时间导数。相容性使 $(K_0u)^Tdy$ 闭，因而在星形坐标域内存在局部标量 $\mathcal h$ 使 $u=J\nabla\mathcal h$。辅助求解同时积分 $\mathsf R'=Dv\mathsf R$ 与 $\upsilon'=Dv\upsilon+v_\tau$，初值分别为 $I$、$0$。二阶敏感度 $\mathsf S=D_y^2\rho$ 和 $\mathsf W=D_y\upsilon$ 满足
\begin{equation}
 \begin{aligned}
 \mathsf S'&=Dv\mathsf S+D^2v[\mathsf R,\mathsf R],\\
 \mathsf W'&=Dv\mathsf W+D^2v[\mathsf R,\upsilon]+Dv_\tau\mathsf R,\\
 Du&=\mathsf R^{-1}\{D\widetilde f_a\mathsf R-\mathsf W-\mathsf S[\cdot,u]\}.
 \end{aligned}
 \label{eq:th-chart-sensitivities}
\end{equation}
这些方程给出被表示动力学的完整切线。一个充分的域条件是 $\sup_U\|\widetilde K-K_0\|_2\le\delta_K<1$，于是 $\sigma_{\min}(K_{\lambda,\tau})\ge1-\delta_K$。对 $\|\widetilde b\|/(1-\delta_K)$ 的界进一步控制辅助位移。由浮点系数计算的多项式界可作为该多项式的充分估计，但不是整个物理模型的区间证书。

同伦条件与端点横截性承担不同作用。前者保证整个辅助区间内可构造坐标图；后者保证所选生成坐标能够参数化端点图。即使坐标图准确，也不能保证所有端点投影在较长物理区间内始终正则。因此，数值实现分别检查坐标图匹配、逆复合、坐标图时间导数和隐式切线。辅助求解仅使用缓存一形式及其导数，不通过积分无法获得的精确物理流来构造坐标图。

\subsection{端点相容性与固定坐标生成}
% Source: stage5/proofs/corrected_mathematics.md, sections 4--5, A1--A4 and B1--B4; stage4/audit/birkhoff_full_audit.md, sections 5--6; original1025 p0862--p0888 and p1179--p1213.
\paragraph{定理 A（局部端点生成函数）.}\label{thm:endpoint}
设 $\Omega_s$ 和 $\Omega_\tau$ 分别为两个端点域上的闭非退化二形式。设局部微分同胚 $\Phi$ 满足 $\Phi^*\Omega_\tau=\Omega_s$，局部可逆的端点变换 $\mathcal E:(z_+,z_-)\mapsto(\widehat w,w)$ 满足
\begin{equation}
 \mathcal E^*d(\widehat w^Tdw)=\Omega_\tau(z_+)-\Omega_s(z_-).
 \label{eq:th-endpoint-pullback}
\end{equation}
记 $D\mathcal E=\left(\begin{smallmatrix}A&B_e\\C&D_e\end{smallmatrix}\right)$、$M_\Phi=D\Phi$。若 $CM_\Phi+D_e$ 非奇异，则在充分小的星形 $w$ 域上存在标量 $\phi$，使
\begin{equation}
 \begin{aligned}
 \widehat w&=\nabla_w\phi,\qquad N=\phi_{ww},\\
 N&=(AM_\Phi+B_e)(CM_\Phi+D_e)^{-1}=N^T.
\end{aligned}
 \label{eq:th-endpoint-hessian}
\end{equation}
反之，由 $(z_+,z_-)=\mathcal E^{-1}(\nabla\phi,w)$ 重构端点时，只要负端点投影的导数非奇异，就得到保持二形式的映射。

事实上，将 \eqref{eq:th-endpoint-pullback} 拉回到 $(\Phi(z_-),z_-)$ 使 $\widehat w^Tdw$ 成为闭一形式。横截性使 $w$ 可作坐标，局部 Poincar'e 引理给出 $\phi$。分别对两个端点变量微分即可得到其 Hessian。重构使用联合端点逆；当下式所需分块可逆时，
\begin{equation}
 M_\Phi=(A-NC)^{-1}(ND_e-B_e).
 \label{eq:th-endpoint-reconstruction}
\end{equation}
无需对 $N$ 求逆。在正则中心坐标的恒等映射处，$N=0$，而所需投影仍然可逆。

初始生成函数也必须由此几何结构确定。反解 $w_0(z)=\mathcal E_2(z,z,s,s)$，并令 $c_0(w;s)=\mathcal E_1(z(w),z(w),s,s)$。闭性给出归一化势函数
\begin{equation}
 \phi_0(w;s)=\int_0^1 c_0(w_*+u(w-w_*);s)^T(w-w_*)\,du .
 \label{eq:th-initial-gauge}
\end{equation}
一般而言 $c_0$ 非零。只有指定随时间变化的加性规范后，势函数才唯一。初始横截性仅保证一个局部时间区间；继续延拓要求联合逆和两个投影保持正则。

例如，将正则中心变量改为 $\widehat w=J(y_+-y_-)+\nabla\psi(w)$ 时，因为 $d(\nabla\psi^Tdw)=0$，端点二形式不变。在对角线上得到 $c_0=\nabla\psi$，而不是零，其归一化初始势为 $\psi(w)-\psi(w_*)$。这个例子区分了端点规范的选择自由度和独立的横截条件。采样端点 Jacobian 非奇异仅说明该采样点正则；统一延拓需要邻域界，防止其最小奇异值趋于零。

\paragraph{定理 B（固定 $w$ 的演化与截断）.}\label{thm:fixed-w}
在共同局部解析域上满足定理 A 的假设，令 $g=\nabla_w\phi$、$N=D_wg$。以下量都通过 $(g,w,\tau,s)$ 处的完整端点逆评价：
\begin{equation}
 \begin{aligned}
 \left.\partial_\tau g\right|_w
 &=(A-NC)f(z_+,\tau)\\
 &\quad+\mathcal E_{1,\tau}-N\mathcal E_{2,\tau}\\
 &=:\Psi(g,w,N,\tau;s).
\end{aligned}
 \label{eq:th-fixed-w}
\end{equation}
这由 $\dot{\widehat w}=A f+\mathcal E_{1,\tau}$、$\dot w=Cf+\mathcal E_{2,\tau}$ 和 $\dot{\widehat w}=N\dot w+\partial_\tau g|_w$ 直接得到。对 $g=\sum_{j\ge0}\theta^jc_j$，可计算递推为
\begin{equation}
 \begin{aligned}
 c_{m+1}&=\frac1{m+1}[\theta^m]\Psi\Bigl(
 \sum_{j=0}^m\theta^jc_j,w,\\
 &\hspace{2em}\sum_{j=0}^m\theta^jD_wc_j,s+\theta;s\Bigr).
\end{aligned}
 \label{eq:th-fixed-recursion}
\end{equation}
特别地，$2c_2=(\Psi_\tau+\Psi_g c_1+\Psi_N[D_wc_1])_0$。精确 $N$ 的对称性给出每个 $D_wc_j$ 的对称性。因此每个有限截断仍是标量梯度，只要投影保持正则，就重构出保持二形式的局部映射。局部流的解析依赖和解析端点逆保证局部展开收敛；有限光滑性只能支持相应的有限 Taylor 余项。精度还需要统一的逆映射界及单步稳定性。

对正则端点 $y_\pm$，取 $w=(y_++y_-)/2$、$\widehat w=J(y_+-y_-)$，联合逆为 $y_\pm=w\mp J\nabla\phi/2$。固定标量的时间规范后，\eqref{eq:th-fixed-w} 化为
\begin{equation}
 \begin{aligned}
 \partial_\tau\phi&=-\mathcal h\left(w-\tfrac12J\nabla\phi,\tau\right),\\
 \phi(w,s;s)&=0.
\end{aligned}
 \label{eq:th-central-pde}
\end{equation}
从此标量方程提取系数会自动保持梯度可积性。显式时间依赖通常产生偶次项。已实现的 24 维输运分支是具有时变坐标图的自治光滑场；这些方程不代表已实现任意非自治载荷或接触切换。

具体地，写成 $\phi=\sum_{j\ge1}\theta^j\phi_j$，得到
\begin{equation}
 \begin{aligned}
 \phi_{m+1}&=-\frac1{m+1}[\theta^m]\mathcal h\Bigl(\\
 &\quad w-\tfrac12J\nabla\sum_{j=1}^{m}\theta^j\phi_j,s+\theta\Bigr).
\end{aligned}
 \label{eq:th-scalar-recursion}
\end{equation}
对自治 Hamilton 函数，前几个系数为 $\phi_1=-h$、$\phi_2=0$ 和 $\phi_3=-f^Th''f/24$。这些系数依赖给定 Hamilton 函数的导数，无需未知精确流的值。标量递推还说明，为何不能将自治高阶中心公式直接用于时变正则坐标图：坐标图的显式时间依赖也进入同一系数提取过程。对于所选分支，仍须确认其实现及导数计算成本。

\subsection{力学中点法与非可分四阶映射}
% Source: stage5/proofs/corrected_mathematics.md, section 6, C30--C36; original1025 p0910/m0400 through p1003/m0442; stage3/program/rom_integrators.py, midpoint.
为避免与状态混淆，将原关系式（30）—（34）的端点变量写成 $\widehat\xi=P_+-P_-$、$\widehat\pi=Q_+-Q_-$、$\xi=\overline Q$ 和 $\pi=-\overline P$。其联合逆为 $Q_\pm=\xi\pm\widehat\pi/2$、$P_\pm=-\pi\pm\widehat\xi/2$。对于正则保守场，标量 $-\eta h(\xi,-\pi,\tau_m)$ 生成完整时间中点法。一般有阻尼的物理中点法不会仅因采用这些端点记号就获得标量生成函数。

这一障碍在单自由度中已可见。对于 $\dot Q=P$、$\dot P=-Q-\gamma P$，在 $(\xi,\pi)$ 变量中的候选端点梯度为 $(-\xi+\gamma\pi,-\pi)$。其混合导数分别为 $\gamma$ 和零。阻尼非零时，不存在具有该梯度的局部标量。因此只修正某个阻尼符号不能恢复缺失的相容性。必须先使时变二形式和坐标图表示所选物理场，才能在这些适配坐标中应用端点定理。

对完整变质量系统，修正后的力学方程为
\begin{equation}
 \begin{aligned}
 \delta Q&=\eta M(\overline Q)^{-1}\overline P,\\
 \delta P_i&=\eta\Bigl[-\tfrac12\overline P^T
                  \partial_iM^{-1}(\overline Q)\overline P\\
 &\hspace{3em}-V_{,i}(\overline Q)+r_i(\overline z,\tau_m)\Bigr].
\end{aligned}
 \label{eq:th-mechanical-midpoint}
\end{equation}
质量、几何惯性和反馈均在同一未知中点求值。固定求解器预条件矩阵不改变此残差。仅在质量、刚度 $K_{\mathrm{lin}}$、阻尼 $C_d$ 为常矩阵，且中点力 $F_m$ 预先给定时，消元得到
\begin{equation}
 \begin{aligned}
 E&=\tfrac{\eta^2}{4}K_{\mathrm{lin}}M^{-1}+\tfrac\eta2 C_dM^{-1},\\
 (I+E)P_+&=(I-E)P_- -\eta K_{\mathrm{lin}}Q_-+\eta F_m,\\
 Q_+&=Q_-+\tfrac\eta2 M^{-1}(P_++P_-).
 \end{aligned}
 \label{eq:th-leftsolve}
\end{equation}
这是要求 $I+E$ 非奇异的左线性求解。交换矩阵次序或使用标量除法都会改变方法；依赖状态的中点反馈仍是隐式的。

对光滑的实际闭环场 $f_c$，残差为 $z_+-z_- -\eta f_c(\overline z,\tau_m)$，其导数为 $I-\eta Df_c(\overline z,\tau_m)/2$。因此反馈律必须在未知中点评价，且其状态导数必须进入切线。在所选分支上，残差导数非奇异给出局部可解性。统一压缩映射论证还要求所选域被映回自身；某个计算状态处良好的条件数不能替代这一要求。

% Source: stage5/proofs/integrators_full_coil.md, sections 3--4; stage3/source_inputs/v10/v10_framework_20261028/english/sections/supp_structure.tex, central4 and exact-primitive subsections; stage3/program/rom_integrators.py, central_scalar and central4.
\paragraph{命题（Central4 及其原函数）.}\label{thm:central4}
对自治光滑 Hamilton 函数 $h$，定义 $f=J\nabla h$ 和
\begin{equation}
 \begin{aligned}
 S_\eta&=\eta h+\frac{\eta^3}{24}f^Th''f,\\
 z_+-z_-&=J\nabla S_\eta\bigl((z_++z_-)/2\bigr).
\end{aligned}
 \label{eq:th-central4}
\end{equation}
在光滑单值隐式分支上，若所需导数有界且满足统一局部稳定性，则该对称方法为四阶并恰当辛。不要求动能与势能可分。

为说明修正项，令 $z_c$ 为精确时间中点状态，$m$ 为端点均值。Taylor 展开给出 $m=z_c+\eta^2f'f(z_c)/8+O(\eta^4)$。从端点差中消去 $z_c$，得到
\begin{equation}
 \begin{aligned}
 \delta z&=\eta f(m)+\frac{\eta^3}{24}
          \{f''[f,f]-2(f')^2f\}(m)\\
 &\quad+O(\eta^5).
\end{aligned}
 \label{eq:th-central-expansion}
\end{equation}
记 $A_h=h''$，对 $q_h=f^TA_hf$ 求导得 $\nabla q_h=h^{(3)}[f,f]-2A_hJA_hf$，其 $J$ 倍恰好就是三次修正。因此位形相关质量的混合导数是公式的一部分。这里的正号属于前向中心生成函数修正。

对 $\lambda_0=P^TdQ$，更强的恰当辛恒等式具有明确原函数：
\begin{equation}
 \begin{aligned}
 \Phi_\eta^*\lambda_0-\lambda_0
 &=\overline P^Td(\delta Q)+\delta P^Td\overline Q\\
 &=d\{\overline P^T\delta Q-S_\eta(m)\}.
\end{aligned}
 \label{eq:th-primitive}
\end{equation}
它要求生成函数和隐式分支单值。在具有非零周期的角变量域上，仅有闭性不能给出全局原函数。Cayley 切线 $(I-JS_\eta''/2)^{-1}(I+JS_\eta''/2)$ 描述理想隐式映射；有限迭代残差须单独计入。

统一动量阻尼具有另一种精确结构：对常数 $\gamma$ 与 $r=-\gamma P$，取 $\Pi=e^{\gamma\tau}P$、$\mathcal h=e^{\gamma\tau}h(Q,e^{-\gamma\tau}\Pi)$，得到正则时变系统。其中点更新为 $S_{\eta/2}\circ\mathrm{MP}_{h,\eta}\circ S_{\eta/2}$，这里 $S_a(Q,P)=(Q,e^{-\gamma a}P)$。它在 $D\Phi^TK_0D\Phi$ 中保持因子 $e^{-\gamma\eta}$，并且为二阶。此特例不能将该因子推广到一般阻尼或磁反馈。

\subsection{有条件的有限时间精度}
% Source: stage5/proofs/integrators_full_coil.md, section 5; stage5/proofs/birkhoff_transport.md, section 5; stage3/source_inputs/v10/v10_framework_20261028/english/sections/supp_structure.tex, implicit-solvability subsection.
\paragraph{命题（分离的有限时间缺陷）.}\label{thm:finite-time}
假设在 $[0,T]$ 上存在共同光滑域，坐标图及其逆的导数统一有界，$\|\widetilde f_a-f\|\le\epsilon_f$，且隐式分支稳定。设中点不动点映射将该域内的一个闭凸子集映回自身，并满足 $\eta L/2<1$。若 $\varepsilon_k$ 为代数增量残差，则相减精确与数值中点方程得到
\begin{equation}
 (1-\eta L/2)e_{k+1}\le(1+\eta L/2)e_k+C\eta^3+\|\varepsilon_k\|.
 \label{eq:th-residual-recursion}
\end{equation}
因此积分误差被 $C_T(e_0+\eta^2+\sum_k\|\varepsilon_k\|)$ 控制。Central4 在相应导数与稳定性假设下将 $\eta^2$ 替换为 $\eta^4$。若 $\delta_k$ 界定数值坐标图操作引入的附加状态缺陷，则物理状态误差满足
\begin{equation}
 \begin{aligned}
 \max_{k\eta\le T}\|z_k-z(k\eta)\|
 &\le C_T\bigl(e_0+\eta^p\bigr)\\
 &\quad+C_T\sum_k(\|\varepsilon_k\|+\delta_k)\\
 &\quad+\epsilon_f\frac{e^{LT}-1}{L}.
\end{aligned}
 \label{eq:th-error-bound}
\end{equation}
其中坐标图中点法取 $p=2$，自治保守 Central4 取 $p=4$；当 $L=0$ 时，末项商的取值为 $T$。统一单步残差 $\varepsilon_\eta$ 贡献 $O(\varepsilon_\eta/\eta)$。这同时解释了表示误差平台和固定求解容差下渐近阶数的丧失。逐点坐标图诊断和方向导数检验支持实现，但不认证所有统一假设。此定理仅涉及有限时间，既不提供全局坐标图构造，也不构成新的 KAM 证书。

此误差估计还决定了参考轨迹的比较方式。原场的参考轨迹测量总差异，被表示场的参考轨迹用于分离积分精度。将后者与坐标图更新比较，可测量数值坐标图与中点法的合并贡献。每项比较必须采用相同初态和时间区间。向其中一条参考轨迹收敛，不能替代与另一条参考轨迹的比较。
